\chapter{Le modèle : un moteur de jeu autonome}
\label{chap-modele}

Le modèle est une bibliothèque de \num{1343} lignes, écrite en C11~\cite{iso9899}, qui connaît
les règles de la bataille navale et rien d'autre. Ce chapitre en présente les structures de données, puis les
traitements, dans l'ordre de la chaîne de modules de la \cref{fig-modules}.

\section{Représenter un plateau}

Dans la version~1, une case était décrite par une structure mêlant l'état du jeu et
l'affichage (\cref{lst-v1-cell}). Tester si un bateau était coulé supposait d'avoir construit
des boutons.

\begin{lstlisting}[caption={Version~1 : la case connaît son bouton.}, label={lst-v1-cell}]
typedef struct {
    GtkWidget *button;  // Button to display the cell
    gboolean is_ship;   // Flag indicating whether there is a ship in this box
    gboolean is_hit;
    gboolean is_shot;
} Cell;
\end{lstlisting}

La version~2 sépare deux natures d'information que la première confondait : le
\textbf{secret} --- où sont les bateaux --- et ce qui est \textbf{public} --- ce que les
tirs ont révélé (\cref{lst-board}).

\begin{lstlisting}[caption={Le plateau d'un camp (\fichier{include/battleship/board.h}).},
                   label={lst-board}]
typedef struct bs_ship {
    bs_coord origin;            /* pont le plus en haut à gauche */
    bs_orientation orientation; /* direction des ponts suivants   */
    uint8_t length;             /* nombre de ponts, de 1 à 4      */
    uint8_t hits;               /* ponts déjà touchés             */
} bs_ship;

typedef struct bs_board {
    bs_ship ships[BS_FLEET_SIZE];
    uint8_t ship_count;            /* bateaux placés                    */
    uint8_t ships_afloat;          /* bateaux non coulés                */
    uint8_t decks_afloat;          /* ponts non touchés                 */
    int8_t  ship_at[BS_CELL_COUNT]; /* secret : indice du bateau, ou -1 */
    uint8_t state[BS_CELL_COUNT];   /* public : ce que la case a révélé */
} bs_board;
\end{lstlisting}

L'état public d'une case prend l'une des cinq valeurs du \cref{tab-etats}. La cinquième,
\code{EXCLUDED}, n'existait pas dans la version~1 comme donnée : elle n'apparaissait que
sous la forme d'un drapeau posé sur un bouton.

\begin{table}[ht]
  \centering
  \begin{tabularx}{\textwidth}{@{}lX@{}}
    \toprule
    \textbf{État} & \textbf{Signification pour l'attaquant} \\
    \midrule
    \code{BS\_CELL\_UNKNOWN}  & Jamais visée, rien n'est connu. \\
    \code{BS\_CELL\_MISS}     & Visée : de l'eau. \\
    \code{BS\_CELL\_HIT}      & Visée : pont d'un bateau encore à flot. \\
    \code{BS\_CELL\_SUNK}     & Visée : pont d'un bateau coulé. \\
    \code{BS\_CELL\_EXCLUDED} & Non visée, mais \emph{prouvée} vide par la règle de
                                non-contact. \\
    \bottomrule
  \end{tabularx}
  \caption{États publics d'une case.}
  \label{tab-etats}
\end{table}

\begin{retenir}[Invariants du plateau]
Toutes les fonctions du module préservent trois propriétés, que les tests de propriétés
revérifient après chaque tir de \num{2250} parties (\cref{sec-proprietes}) :
\begin{itemize}
  \item les bateaux ne se chevauchent pas, ne se touchent pas et ne sortent pas du plateau ;
  \item \code{decks\_afloat} est égal au nombre de ponts non encore touchés ;
  \item une case est \code{SUNK} si et seulement si elle appartient à un bateau dont tous
        les ponts sont touchés.
\end{itemize}
\end{retenir}

\subsection{Aucune allocation dynamique}

Les dimensions du jeu sont fixes : tout tient dans des tableaux de taille connue à la
compilation. Une partie entière, \code{bs\_game}, occupe 800~octets. Le moteur n'appelle
jamais \code{malloc}. Il n'y a donc ni échec d'allocation à traiter, ni fuite possible, et
copier une partie se réduit à une affectation. C'est ce qui rend bon marché l'exécution de
millions de parties simulées.

\section{Placer une flotte au hasard}

\subsection{L'algorithme de la version~1 : tirer jusqu'à ce que ça rentre}

La première version tirait une position et une orientation au hasard, vérifiait si le
bateau tenait, et recommençait sinon, jusqu'à un million de tentatives. La méthode
fonctionne, mais sa durée n'est pas bornée de façon utile et son test de validité comportait
une erreur de borne, masquée par le choix des intervalles de tirage.

\subsection{L'algorithme de la version~2 : énumérer, puis choisir}

Pour chaque bateau, du plus long au plus court, on dresse la liste de \emph{toutes} les
positions licites, puis on en choisit une uniformément (\cref{alg-placement}). Une grille
de 100 cases et deux orientations donnent au plus 200 candidats : l'énumération est
immédiate, et l'algorithme ne dépend plus de la chance pour terminer.

\begin{algorithm}[ht]
  \caption{Placement aléatoire d'une flotte}
  \label{alg-placement}
  \KwIn{un générateur aléatoire $g$}
  \KwOut{un plateau portant les dix bateaux, sans contact}
  \Repeat{la flotte est complète}{
    vider le plateau\;
    \ForEach{longueur $\ell$ de la flotte, dans l'ordre $4, 3, 3, 2, 2, 2, 1, 1, 1, 1$}{
      $P \gets$ ensemble des couples (origine, orientation) où un bateau de longueur $\ell$
      tient sans toucher un bateau déjà placé\;
      \lIf{$P = \emptyset$}{abandonner cette tentative \Com*[f]{impasse}}
      choisir $p \in P$ uniformément avec $g$ et placer le bateau en $p$\;
    }
  }
\end{algorithm}

Le test \enquote{tient sans toucher} est une simple inspection : on parcourt le rectangle
occupé par le bateau, agrandi d'une case dans toutes les directions, et l'on refuse la
position dès qu'une case de ce rectangle porte déjà un bateau. Un bateau d'un seul pont est
identique dans les deux orientations ; il n'est compté qu'une fois, pour ne pas fausser le
tirage.

Une impasse --- plus aucune place pour un bateau restant --- est théoriquement possible.
Elle ne s'est produite dans aucune des flottes tirées par les tests, soit plusieurs
dizaines de milliers ; la boucle externe, bornée à mille tentatives, ne sert qu'à transformer un
\enquote{ne peut pas arriver} en erreur signalée plutôt qu'en boucle sans fin.

\section{Tirer, et en tirer des conclusions}
\label{sec-deductions}

Un tir met à jour l'historique public puis les \emph{déductions} que la règle de
non-contact autorise (\cref{fig-deductions}) :

\begin{itemize}
  \item \textbf{touché} --- les quatre voisines diagonales de la case sont de l'eau, car
        un bateau est droit et aucun autre ne peut le frôler ;
  \item \textbf{coulé} --- toute la ceinture du bateau est de l'eau.
\end{itemize}

\begin{figure}[ht]
  \centering
  \begin{tikzpicture}[scale=0.62]
    % ----------------------------------------------------- (a) après un coup au but
    \begin{scope}
      \plateau{7}
      \fill[ocean!25] (2,3) rectangle (5,4);
      \draw[ocean, dashed, line width=0.6pt, rounded corners=1.5pt] (2.1,3.1) rectangle (4.9,3.9);
      \fill[corail] (3.5,3.5) circle (0.24);
      \foreach \x/\y in {2/2, 4/2, 2/4, 4/4} {\drapeau{\x}{\y}}
      \node[note] at (3.5,-0.8) {(a) touché :\\4 cases exclues};
    \end{scope}
    % ------------------------------------------------------ (b) après le naufrage
    \begin{scope}[xshift=9.5cm]
      \plateau{7}
      \fill[marine, rounded corners=1.5pt] (2.12,3.12) rectangle (4.88,3.88);
      \foreach \x in {2,3,4} \fill[corail] (\x+0.5,3.5) circle (0.2);
      \foreach \x in {1,...,5} { \drapeau{\x}{2} \drapeau{\x}{4} }
      \drapeau{1}{3} \drapeau{5}{3}
      \node[note] at (3.5,-0.8) {(b) coulé :\\12 cases exclues};
    \end{scope}
  \end{tikzpicture}
  \caption{Déductions tirées d'un tir. Les croix marquent les cases que le moteur classe
           \code{EXCLUDED} ; en (a) le reste du bateau, en pointillé, est encore inconnu.}
  \label{fig-deductions}
\end{figure}

Ces déductions sont calculées \emph{par le modèle}, au moment du tir, et non par chaque
interface. Dans la version~1, le marquage des drapeaux était recopié dans chacun des trois
gestionnaires de clic, une fois pour le joueur et une fois pour l'ordinateur : c'était à
lui seul près de la moitié de leur longueur. Il tient désormais en deux fonctions d'une
dizaine de lignes.

\section{Ce que l'attaquant a le droit de savoir}
\label{sec-connaissance}

Pour que l'adversaire artificiel ne puisse pas tricher, il ne suffit pas de lui demander
de ne pas regarder. Le modèle extrait du plateau une structure distincte, qui ne contient
que l'information publique (\cref{lst-knowledge}).

\begin{lstlisting}[caption={La connaissance publique d'un plateau.}, label={lst-knowledge}]
typedef struct bs_knowledge {
    /* état public de chaque case */
    uint8_t state[BS_CELL_COUNT];
    /* nombre de bateaux encore à flot, par longueur */
    uint8_t afloat_by_length[BS_MAX_SHIP_LENGTH + 1];
} bs_knowledge;
\end{lstlisting}

L'adversaire reçoit cette structure et rien d'autre. La position des bateaux ne fait pas
partie de ses arguments : l'exigence~C3 est satisfaite par les types, sans discipline
particulière à respecter.

\section{La partie est un automate}

Le module \code{game} réunit deux plateaux, le niveau de l'adversaire, le score et le tour.
Son comportement est celui de l'automate de la \cref{fig-automate}.

\begin{figure}[ht]
  \centering
  \begin{tikzpicture}[node distance=30mm,
      etat/.style={boite, minimum width=3.0cm, minimum height=1.0cm, rounded corners=5mm}]
    \node[etat, fill=ecume] (h) {Tour du joueur};
    \node[etat, fill=ecume, right=of h] (c) {Tour de l'ordinateur};
    \node[etat, coeur, below=18mm of $(h)!0.5!(c)$] (o) {Partie terminée};
    \node[circle, fill=marine, inner sep=2.2pt, left=26mm of h] (start) {};

    \draw[fleche] (start) -- node[above, note] {initialisation} (h);
    \draw[fleche] (h) to[bend left=14] node[above, note] {manqué} (c);
    \draw[fleche] (c) to[bend left=14] node[below, note] {manqué} (h);
    \draw[fleche] (h) to[loop above, looseness=5] node[above, note] {touché} (h);
    \draw[fleche] (c) to[loop above, looseness=5] node[above, note] {touché} (c);
    \draw[fleche] (h) |- node[pos=0.25, left, note] {flotte adverse\\coulée} (o);
    \draw[fleche] (c) |- node[pos=0.25, right, note] {flotte du joueur\\coulée} (o);
  \end{tikzpicture}
  \caption{L'automate d'une partie. Un coup au but laisse la main au tireur.}
  \label{fig-automate}
\end{figure}

Le principe de l'interface du module est qu'\textbf{un appel effectue un pas, et un seul},
puis rend compte de ce qui a changé (\cref{lst-game-api}). Le moteur ne bloque jamais, ne
temporise jamais et n'appelle personne en retour. Le rythme est l'affaire de l'appelant :
l'interface graphique attend une demi-seconde entre deux tirs de l'ordinateur, le
simulateur n'attend pas du tout.

\begin{lstlisting}[basicstyle=\ttfamily\footnotesize,
                   caption={L'essentiel de l'interface d'une partie
                            (\fichier{include/battleship/game.h}).}, label={lst-game-api}]
bs_status bs_game_init(bs_game *game, bs_difficulty difficulty, uint64_t seed);
bs_status bs_game_reshuffle(bs_game *game);
bs_status bs_game_human_fire(bs_game *game, bs_coord target, bs_turn_report *report);
bs_status bs_game_computer_fire(bs_game *game, bs_turn_report *report);
bs_phase  bs_game_phase(const bs_game *game);
\end{lstlisting}

\subsection{Les erreurs sont des valeurs}

Toute opération qui peut être refusée renvoie un \code{bs\_status} : succès, argument
invalide, position hors du plateau, collision, case déjà jouée, action interdite dans la
phase en cours. Une opération refusée \emph{ne modifie rien} (exigence~C4). Tirer hors de
son tour, par exemple, ne change ni le plateau ni le score ; un test le vérifie
explicitement. Le contrôleur peut ainsi transmettre au modèle n'importe quel clic sans le
filtrer : le modèle est le seul juge.

\section{Un hasard reproductible}
\label{sec-hasard}

\subsection{Le hasard est un argument}

Le moteur n'appelle jamais \code{rand()}. Chaque décision aléatoire puise dans un
générateur, \code{bs\_rng}, que l'appelant possède et initialise par une graine. Deux
conséquences :

\begin{itemize}
  \item une graine rejoue une partie entière --- placement des deux flottes et tous les
        coups de l'ordinateur ; c'est la condition pour qu'un test portant sur une
        stratégie aléatoire soit déterministe ;
  \item le résultat ne dépend pas de la bibliothèque C de la plate-forme, dont le
        \code{rand()} varie d'un système à l'autre.
\end{itemize}

Le générateur est \emph{xoshiro256**}~\cite{blackman2021}, dont l'état de 256~bits est
rempli à partir de la graine par \emph{SplitMix64}~\cite{steele2014}. Ses premières sorties
sont consignées dans les tests : toute différence entre deux plates-formes serait détectée.

\subsection{Tirer un entier sans biais}

La version~1 obtenait un entier de $[0, n[$ par \code{rand() \% n}. L'opération est biaisée
dès que $n$ ne divise pas le nombre de valeurs possibles du générateur. Avec un générateur
sur 3~bits et $n = 3$, par exemple, les huit valeurs $0, \dots, 7$ donnent les restes
$0, 1, 2, 0, 1, 2, 0, 1$ : le reste $2$ sort deux fois sur huit quand les autres sortent
trois fois.

Le module utilise la méthode de Lemire~\cite{lemire2019}, qui multiplie au lieu de diviser
et rejette les quelques tirages responsables du biais (\cref{lst-lemire}). Un test vérifie
l'uniformité sur \num{60000} tirages d'un dé.

\begin{lstlisting}[caption={Entier uniforme dans $[0, \mathit{bound}[$
                            (\fichier{src/core/rng.c}).}, label={lst-lemire}]
uint64_t product = (uint64_t)next_u32(rng) * bound;
uint32_t low = (uint32_t)product;
if (low < bound) {
    const uint32_t threshold = (0U - bound) % bound;
    while (low < threshold) {            /* tirage biaisé : on le rejette */
        product = (uint64_t)next_u32(rng) * bound;
        low = (uint32_t)product;
    }
}
return (uint32_t)(product >> 32);
\end{lstlisting}
